% Unit 063 terjemahan Bahasa Indonesia dari chapter4.tex baris 2500--2745.
% Sumber: Methods of Algebra, Volume 2, commit
% 9a5803ff77dd3257484cb177f851a73770a59dd3 (CC BY 4.0).
% stable-unit-id: o014.aljabr2.chapter4.k-flat-derived-tensor-b
% Cakupan: akhir Bagian 4.12 dan Latihan Bab 4 lengkap (17 latihan dan 10 petunjuk);
% menutup chapter4.tex tepat pada penutup lingkungan Exercises.

% segment-id: o014.li2.chapter4.k-flat-derived-tensor-b.co001
\hypertarget{o014.aljabr2.chapter4.k-flat-derived-tensor-b}{ }
\begin{korolari}\label{prop:otimesL-Tor}
	\index[sym1]{Tor}
	Funktor $\otimesL[R]$ dari Definisi~\ref{def:otimesL}, dengan syarat $A$
	atau $B$ datar sebagai modul-$\Bbbk$, memenuhi isomorfisme kanonik
	\[ \Hm^{-n}\left( X \otimesL[R] Y \right) \simeq \Tor^R_n(X, Y), \quad n \in \Z, \]
	dengan $X$ suatu bimodul-$(A,R)$, $Y$ suatu bimodul-$(R,B)$, dan
	$\Tor^R_n$ seperti dalam Definisi--Proposisi~\ref{def:Tor-classical}.
\end{korolari}
% segment-id: o014.li2.chapter4.k-flat-derived-tensor-b.pf001
\begin{bukti}
	Ruas kiri diperoleh dengan menyusun diagram Proposisi
	\ref{prop:otimesL-enrichment} sepanjang arah
% segment-id: o014.li2.chapter4.k-flat-derived-tensor-b.g001
	\begin{tikzpicture}[xscale=0.3, yscale=0.3, baseline=(O)]
		\draw[->] (0,0) -- (1,0) -- (1, -1);
		\coordinate (O) at (0, -0.7);
	\end{tikzpicture}
	sedangkan ruas kanan diperoleh dengan menyusunnya sepanjang arah
% segment-id: o014.li2.chapter4.k-flat-derived-tensor-b.g002
	\begin{tikzpicture}[xscale=0.3, yscale=0.3, baseline=(O)]
		\draw[->] (0,0) -- (0, -1) -- (1, -1);
		\coordinate (O) at (0, -0.7);
	\end{tikzpicture}.
\end{bukti}

% segment-id: o014.li2.chapter4.k-flat-derived-tensor-b.pr001
\begin{proposisi}[Kendala asosiativitas]\label{prop:otimesL-associativity-constraint}
	\index{hasil kali tensor turunan!kendala asosiativitas (associativity constraint)}
	Misalkan $A$, $B$, $R$, dan $S$ aljabar-$\Bbbk$. Tinjau dua funktor berikut:
	\begin{gather*}
		\cate{D}((A,R)\dcate{Mod}) \times \cate{D}((R,S)\dcate{Mod}) \times \cate{D}((S,B)\dcate{Mod}) \to \cate{D}((A,B)\dcate{Mod}), \\
		(X, Y, Z) \mapsto (X \otimesL[R] Y) \otimesL[S] Z, \\
		(X, Y, Z) \mapsto X \otimesL[R] (Y \otimesL[S] Z).
	\end{gather*}
	Jika $A$ dan $B$ keduanya datar sebagai modul-$\Bbbk$, terdapat
	isomorfisme kanonik
	\[ a(X,Y,Z): (X \otimesL[R] Y) \otimesL[S] Z \rightiso X \otimesL[R] (Y \otimesL[S] Z). \]
\end{proposisi}
% segment-id: o014.li2.chapter4.k-flat-derived-tensor-b.pf002
\begin{bukti}
	Tetapkan $Y$. Berdasarkan Proposisi~\ref{prop:K-flat}, kita boleh
	mengandaikan bahwa $X_R$ dan ${}_R Z$ K-datar. Dengan demikian, semua
	$\otimesL[R]$ dapat dihitung pada tingkat kompleks, dan semuanya kembali
	kepada kendala asosiativitas hasil kali tensor
	\cite[Proposisi 6.5.12]{Li1}.
\end{bukti}

% segment-id: o014.li2.chapter4.k-flat-derived-tensor-b.p001
Dengan asumsi kedataran tersebut, mengambil $\otimesL$ dari lebih dari tiga
objek dalam berbagai urutan memberikan hasil-hasil yang saling isomorfik
melalui kendala asosiativitas. Isomorfisme-isomorfisme ini juga memenuhi
koherensi, yang dapat dirumuskan dengan meniru aksioma segilima untuk kategori
monoidal; lihat \cite[Definisi 3.1.1]{Li1}. Caranya tetap dengan mengambil
	resolusi K-datar dan melakukan verifikasi pada tingkat kompleks. Untuk
menghemat uraian, buku ini hanya menampilkan kasus khusus gelanggang
komutatif.

% segment-id: o014.li2.chapter4.k-flat-derived-tensor-b.e001
\begin{contoh}[$\otimesL$ untuk gelanggang komutatif]\label{eg:commutative-otimesL}
	\index{hasil kali tensor turunan!kendala komutativitas (commutativity constraint)}
	Misalkan $R$ gelanggang komutatif. Dalam definisi
	$\otimesL := \otimesL[R]$, ambil semua gelanggang sama dengan $R$,
	khususnya $\Bbbk=R$. Dengan demikian,
	$(R,R)\dcate{Mod}=R\dcate{Mod}$ dan semua asumsi kedataran sebelumnya
	terpenuhi secara otomatis. Akibatnya,
	$\Hm^{-n}(X \otimesL Y) \simeq \Tor^R_n(X,Y)$ berlaku untuk semua
	modul-$R$ $X$ dan $Y$.
	
	Hal ini menjadikan $\cate{D}(R\dcate{Mod})$ kategori monoidal terhadap
	$\otimesL$ dalam arti \cite[\S 3.1]{Li1}. Objek unitnya adalah $R$
	sendiri, sedangkan kendala asosiativitasnya berasal dari
	$a=(a(X,Y,Z))_{X,Y,Z}$ dalam Proposisi
	\ref{prop:otimesL-associativity-constraint}. Karena $R$ merupakan
	modul-$R$ datar, ia K-datar (Proposisi~\ref{prop:flat-K-flat}); maka
	$R \otimesL (\cdot)$ dan $(\cdot) \otimesL R$ keduanya dapat dihitung
	pada tingkat kompleks, sehingga semua aksioma kategori monoidal mudah
	dibuktikan. Lebih lanjut, $\cate{D}(R\dcate{Mod})$ merupakan kategori
	monoidal simetris dalam arti \cite[Definisi 3.3.1]{Li1}. Struktur
	anyamannya (atau kendala komutativitasnya)
	$c(X,Y):X\otimesL Y\rightiso Y\otimesL X$ berasal dari kendala
	komutativitas pada tingkat kompleks, dan aksioma-aksiomanya juga diperiksa
	pada tingkat kompleks. Mengingat Proposisi~\ref{prop:double-cplx-swap},
	pertukaran superskrip pada kompleks total tidak menimbulkan masalah di sini.
\end{contoh}

% segment-id: o014.li2.chapter4.k-flat-derived-tensor-b.e002
\begin{contoh}[Aljabar $\mathrm{Tor}$]
	Melanjutkan pembahasan Contoh~\ref{eg:commutative-otimesL}, misalkan $R$
	komutatif. Identifikasikan kompleks modul-$R$ $X,Y,Z,W$ dengan citranya
	dalam kategori turunan. Untuk setiap $(p,q)\in\Z^2$, terdapat morfisme
	kanonik
% segment-id: o014.li2.chapter4.k-flat-derived-tensor-b.q001
	\begin{multline*}
		\Hm^{-p}\left( X \otimesL Y \right) \otimes \Hm^{-q}\left( Z \otimesL W \right) \to \Hm^{-p-q}\left( (X \otimesL Y) \otimesL (Z \otimesL W) \right) \\
		\simeq \Hm^{-p-q}\left( (X \otimesL Z) \otimesL (Y \otimesL W) \right) \to \Hm^{-p-q} \left( (X \otimes Z) \otimesL (Y \otimes W) \right),
	\end{multline*}
	dengan bagian pertama merupakan penerapan Proposisi~\ref{prop:bifunctor-cup}
	(ambil $F=\otimes$), bagian kedua memakai kendala asosiativitas dan
	kendala komutativitas pada $(\cate{D}(R\dcate{Mod}),\otimesL)$, dan
	bagian ketiga berasal dari morfisme kanonik yang melekat pada $\otimesL$
	sebagai bifunktor turunan kiri:
% segment-id: o014.li2.chapter4.k-flat-derived-tensor-b.q002
	\[ QX_1 \otimesL QX_2 \to Q(X_1 \otimes X_2), \quad X_1, X_2 \in \Obj(\cate{K}(R\dcate{Mod})), \]
	Untuk ketepatan notasi, di sini funktor lokalisasi
	$Q:\cate{K}(R\dcate{Mod})\to\cate{D}(R\dcate{Mod})$ dituliskan kembali.
% segment-id: o014.li2.chapter4.k-flat-derived-tensor-b.l001
	\begin{itemize}
% segment-id: o014.li2.chapter4.k-flat-derived-tensor-b.i001
		\item Komposisi ketiga morfisme kanonik di atas memberikan ``hasil kali
		luar'' bagi funktor $\Tor$:
		\[ \Tor^R_p(X, Y) \otimes \Tor^R_q(Z, W) \to \Tor^R_{p+q}(X \otimes Z, Y \otimes W). \]
% segment-id: o014.li2.chapter4.k-flat-derived-tensor-b.i002
		\item Ambil $Z=X$ dan $W=Y$ sebagai aljabar-$R$ (terkonsentrasi pada
		derajat $0$). Perkaliannya memberikan homomorfisme modul-$R$
		$X\otimes X\to X$ dan $Y\otimes Y\to Y$. Dengan menyusun hasil kali
		luar dengan homomorfisme-homomorfisme ini, diperoleh
		\[ \Tor^R_p(X, Y) \otimes \Tor^R_q(X, Y) \to \Tor^R_{p+q}(X,Y), \]
		yang disebut ``hasil kali dalam'' bagi funktor $\Tor$. Hal ini menjadikan
		$\bigoplus_{n \in \Z} \Tor^R_n(X, Y)$ aljabar-$R$ bergradasi-$\Z$
		dalam arti \cite[Definisi 7.4.1]{Li1}, yang disebut \emph{aljabar
		$\Tor$}. Aljabar ini hanya memiliki komponen tak nol pada derajat
		nonnegatif, dan komponen derajat $0$-nya memberikan subaljabar
		$\Tor^R_0(X, Y) = X \otimes Y$, yang definisinya dapat dilihat dalam
		\cite[Definisi 7.3.1]{Li1}. Hukum asosiatif yang diperlukan bagi
		perkalian dapat diverifikasi dengan mereduksinya kepada kendala
		asosiativitas $\otimesL$ dan hukum asosiatif perkalian masing-masing
		pada $X$ dan $Y$.
		\index{aljabar Tor@aljabar $\Tor$ ($\Tor$-algebra)}
% segment-id: o014.li2.chapter4.k-flat-derived-tensor-b.i003
		\item Andaikan lebih lanjut bahwa $X$ dan $Y$ merupakan aljabar-$R$
		komutatif. Dalam hal ini, aljabar $\Tor$ juga merupakan aljabar
		bergradasi-$\Z$ ``antikomutatif'' dalam arti
		\cite[Definisi 7.4.3]{Li1}\footnote{Sebutan harfiahnya agak menyesatkan,
		karena sifat ini merupakan perwujudan alami komutativitas. Sebutan
		bergradasi-komutatif lebih tepat, seperti dalam Contoh
		\ref{eg:graded-module}.}: untuk semua $p,q\in\Z$,
		\[ a \in \Tor^R_p(X, Y), \; b \in \Tor^R_q(X, Y) \implies ab = (-1)^{pq} ba. \]
		Karena $X\otimes X\to X$ dan $Y\otimes Y\to Y$ keduanya komutatif,
		dari mana tanda tersebut berasal? Sumbernya adalah definisi hasil kali
		luar. Misalkan $C_1$ dan $C_2$ dua salinan $X\otimesL Y$. Maka $ab$
		dan $ba$ berbeda melalui isomorfisme
		$C_1\otimesL C_2\rightiso C_2\otimesL C_1$, yaitu kendala
		komutativitas $\otimesL$. Jika $C_1$ dan $C_2$ dinyatakan sebagai
		kompleks K-datar, isomorfisme pada tingkat kompleks total ini persis
		isomorfisme $r_{C_1^\bullet \otimes C_2^\bullet}$ dari
		Proposisi~\ref{prop:double-cplx-swap}; pada komponen berderajat $(p,q)$,
		isomorfisme tersebut menghasilkan tanda $(-1)^{pq}$.
	\end{itemize}
\end{contoh}

% segment-id: o014.li2.chapter4.k-flat-derived-tensor-b.p002
Kembali ke aljabar-$\Bbbk$ umum. Untuk membahas adjungsi antara $\otimesL$
dan $\RHom$, masih diperlukan satu lema. Berikut ini, $\RHom_{(A,B)}$
menyatakan $\RHom$ pada $(A,B)\dcate{Mod}$; notasi bagi kategori bimodul
(atau modul kiri maupun modul kanan) jenis lainnya digunakan secara serupa.

% segment-id: o014.li2.chapter4.k-flat-derived-tensor-b.le001
\begin{lema}\label{prop:RHom-enrichment}
	Tetapkan aljabar-$\Bbbk$ $A$, $B$, dan $R$. Andaikan $B$ datar sebagai
	modul-$\Bbbk$. Maka funktor pelupa dari $(A,B)\dcate{Mod}$ ke
	$A\dcate{Mod}$ mempertahankan kompleks K-injektif, dan $\RHom_A$ secara
	kanonik dapat diangkat menjadi funktor
	\[ \cate{D}\left((A,R)\dcate{Mod}\right)^{\opp} \times \cate{D}\left((A,B)\dcate{Mod}\right) \to \cate{D}\left((R,B)\dcate{Mod}\right), \]
	dengan kata lain, terdapat diagram 2-sel kanonik
% segment-id: o014.li2.chapter4.k-flat-derived-tensor-b.g003
	\[\begin{tikzcd}
		\cate{D}\left((A,R)\dcate{Mod}\right)^{\opp} \times \cate{D}\left((A,B)\dcate{Mod}\right) \arrow[r] \arrow[d, "\text{pelupa}"'] & \cate{D}\left((R,B)\dcate{Mod}\right) \arrow[d, "\text{pelupa}"] \arrow[Rightarrow, ld, "\sim", sloped] \\
		\cate{D}(A\dcate{Mod})^{\opp} \times \cate{D}(A\dcate{Mod}) \arrow[r, "{\RHom_A}"'] & \cate{D}(\cate{Ab}) .
	\end{tikzcd}\]

	Demikian pula, jika $A$ datar sebagai modul-$\Bbbk$, $\RHom_B$ mempunyai
	pengangkatan serupa.
\end{lema}
% segment-id: o014.li2.chapter4.k-flat-derived-tensor-b.pf003
\begin{bukti}
	Karena $B$ modul-$\Bbbk$ datar, funktor $Y\mapsto{}_A Y$ mempunyai adjoin
	kiri eksak $(\cdot)\dotimes{\Bbbk}B$, sehingga ia mempertahankan kompleks
	K-injektif.
	
	Selanjutnya, ambil
	$X \in \Obj(\cate{K}((A, R)\dcate{Mod}))$ dan
	$Y \in \Obj(\cate{K}(A, B)\dcate{Mod})$, dengan $Y$ K-injektif.
	Berdasarkan langkah sebelumnya,
	$\Hom^\bullet\left({}_A X,{}_A Y\right)$ menentukan
	$\RHom_A({}_A X,{}_A Y)$, tetapi sekaligus merupakan kompleks
	bimodul-$(R,B)$. Ini menyelesaikan proses pengangkatan tersebut.
\end{bukti}

% segment-id: o014.li2.chapter4.k-flat-derived-tensor-b.th001
\begin{teorema}[Adjungsi antara $\RHom$ dan $\otimesL$]\label{prop:RHom-otimesL-adjunction}
	\index{funktor $\RHom$}
	\index{hasil kali tensor turunan}
	Misalkan $A$, $B$, dan $R$ aljabar-$\Bbbk$. Di bawah ini, $X$, $Y$, dan
	$Z$ masing-masing menyatakan objek
	$\cate{D}((A,R)\dcate{Mod})$, $\cate{D}((R,B)\dcate{Mod})$, dan
	$\cate{D}((A,B)\dcate{Mod})$. Andaikan $B$ datar sebagai
	modul-$\Bbbk$. Maka terdapat isomorfisme kanonik dalam
	$\cate{D}(\Bbbk\dcate{Mod})$
% segment-id: o014.li2.chapter4.k-flat-derived-tensor-b.q003
	\begin{equation*}
		\RHom_{(A,B)} \left( X \otimesL[R] Y, Z \right) \rightiso \RHom_{(R,B)}\left( Y, \RHom_A \left( {}_A X, {}_A Z \right) \right).
	\end{equation*}

	Jika $A$ yang diasumsikan datar sebagai modul-$\Bbbk$, terdapat
	isomorfisme kanonik dalam $\cate{D}(\Bbbk\dcate{Mod})$
% segment-id: o014.li2.chapter4.k-flat-derived-tensor-b.q004
	\begin{equation*}
		\RHom_{(A,B)}\left( X \otimesL[R] Y, Z \right) \rightiso \RHom_{(A,R)}\left( X, \RHom_B\left( Y_B , Z_B \right) \right).
	\end{equation*}
\end{teorema}
% segment-id: o014.li2.chapter4.k-flat-derived-tensor-b.pf004
\begin{bukti}
	Tinjau persamaan pertama. Andaikan $Y$ K-projektif dan $Z$ K-injektif.
	Lema~\ref{prop:K-projective-flat} menjamin bahwa ${}_R Y$ K-datar,
	sedangkan Lema~\ref{prop:RHom-enrichment} menjamin bahwa ${}_A Z$ tetap
	K-injektif. Semua $\RHom$ dan $\otimesL[R]$ yang ditampilkan dapat
	ditentukan pada tingkat kompleks. Karena itu, isomorfisme yang dicari
	mengikuti dari Proposisi~\ref{prop:tensor-Hom-cplx}.
\end{bukti}

% segment-id: o014.li2.chapter4.k-flat-derived-tensor-b.e003
\begin{contoh}[Adjungsi untuk perubahan gelanggang]\label{eg:change-of-ring-adjunction}
	Dalam situasi Contoh~\ref{eg:change-of-rings} (ambil homomorfisme
	aljabar-$\Bbbk$ $R\to S$, $A=S$, $B=\Bbbk$, dan $X=S$), persamaan pertama
	dalam Teorema~\ref{prop:RHom-otimesL-adjunction} menjadi isomorfisme kanonik
% segment-id: o014.li2.chapter4.k-flat-derived-tensor-b.q005
	\begin{gather*}
		\RHom_S\left( \mathrm{L} {}_{R \to S} P(Y), Z \right) \simeq \RHom_R\left( Y, {}_R Z \right), \\
		Y \in \Obj(\cate{D}(R\dcate{Mod})), \quad Z \in \Obj(\cate{D}(S\dcate{Mod})).
	\end{gather*}
	Untuk itu, satu-satunya hal yang masih perlu dijelaskan ialah
	$\RHom_S({}_S S,{}_S Z)\simeq{}_R Z$, yang langsung diperoleh dari
	isomorfisme pada tingkat kompleks
% segment-id: o014.li2.chapter4.k-flat-derived-tensor-b.q006
	\[ \Hom^\bullet\left({}_S S, \cdot\right) \simeq {}_R (\cdot) : \cate{C}(S\dcate{Mod}) \to \cate{C}(R\dcate{Mod}). \]
	
	Jika selanjutnya $R$ dan $S$ keduanya gelanggang komutatif dan
	$\Bbbk:=R$, maka $\RHom_S$ (atau $\RHom_R$) bernilai dalam
	$\cate{D}(S\dcate{Mod})$ (atau $\cate{D}(R\dcate{Mod})$). Isomorfisme
	adjungsi tersebut selanjutnya dapat dituliskan sebagai isomorfisme kanonik
	dalam $\cate{D}(R\dcate{Mod})$:
% segment-id: o014.li2.chapter4.k-flat-derived-tensor-b.q007
	\[ {}_R \RHom_S\left( \mathrm{L} {}_{R \to S} P(Y), Z \right) \simeq \RHom_R\left( Y, {}_R Z \right) . \]
	Semuanya dapat diverifikasi pada tingkat kompleks dengan mengambil
	resolusi K-injektif dan K-projektif. Kasus modul kanan serupa.
\end{contoh}


% segment-id: o014.li2.chapter4.k-flat-derived-tensor-b.x001
\begin{Exercises}
	% Reference: [KS06, Exercise 10.10]
% segment-id: o014.li2.chapter4.k-flat-derived-tensor-b.ex001
	\item Misalkan $(\mathcal{D}, T, \mathcal{H})$ kategori prabertriangulasi
	(atau kategori bertriangulasi). Definisikan keluarga segitiga baru
	$\mathcal{H}^-$ melalui
% segment-id: o014.li2.chapter4.k-flat-derived-tensor-b.q008
	\[ [X \xrightarrow{f} Y \xrightarrow{g} Z \xrightarrow{h} TX] \in \mathcal{H} \iff [X \xrightarrow{f} Y \xrightarrow{g} Z \xrightarrow{-h} TX] \in \mathcal{H}^- . \]
% segment-id: o014.li2.chapter4.k-flat-derived-tensor-b.l002
	\begin{enumerate}[(i)]
% segment-id: o014.li2.chapter4.k-flat-derived-tensor-b.i004
		\item Buktikan bahwa $(\mathcal{D},T,\mathcal{H}^-)$ juga merupakan
		kategori prabertriangulasi (atau kategori bertriangulasi).
% segment-id: o014.li2.chapter4.k-flat-derived-tensor-b.i005
		\item Buktikan bahwa $(\mathcal{D},T,\mathcal{H})$ dan
		$(\mathcal{D},T,\mathcal{H}^-)$ saling ekuivalen sebagai kategori
		prabertriangulasi.
	\end{enumerate}

% segment-id: o014.li2.chapter4.k-flat-derived-tensor-b.ex002
	\item Misalkan $(\mathcal{D},T,\mathcal{H})$ kategori prabertriangulasi.
	Buktikan bahwa setiap monomorfisme (atau epimorfisme) $u:X\to Y$ dalam
	$\mathcal{D}$ mempunyai invers kiri (atau invers kanan). Gunakan hal ini
	untuk menunjukkan bahwa kategori prabertriangulasi yang sekaligus kategori
	abelian harus terbelah (Definisi~\ref{def:semisimple}).
	
% segment-id: o014.li2.chapter4.k-flat-derived-tensor-b.hn001
	\begin{hint}
		Tempatkan monomorfisme $u$ dalam segitiga terbedakan
		$X\xrightarrow{u}Y\xrightarrow{v}Z\xrightarrow{w}TX$. Dari
		$u\circ T^{-1}w=0$, simpulkan $w=0$, lalu terapkan
		Proposisi~\ref{prop:dt-sum-1}.
	\end{hint}

% segment-id: o014.li2.chapter4.k-flat-derived-tensor-b.ex003
	\item Terapkan latihan sebelumnya untuk memberikan contoh kategori abelian
	$\mathcal{A}$ sedemikian sehingga $\cate{K}(\mathcal{A})$ dan
	$\cate{D}(\mathcal{A})$ keduanya bukan kategori abelian.

% segment-id: o014.li2.chapter4.k-flat-derived-tensor-b.hn002
	\begin{hint}
		Tinjau morfisme $f:\Z/p^2\Z\twoheadrightarrow\Z/p\Z$ dalam
		$\mathcal{A}=\cate{Ab}$, dengan $p$ prima. Jika
		$\cate{K}(\mathcal{A})$ kategori abelian, tinjau faktorisasi epi--mono
		$\cate{K}f$ dan $\Hm^0(\cate{K}f)$, kemudian uraikan keduanya sebagai
		proyeksi ke suku jumlah langsung dan morfisme inklusi untuk memperoleh
		kontradiksi.
	\end{hint}

% segment-id: o014.li2.chapter4.k-flat-derived-tensor-b.ex004
	\item Misalkan $\mathcal{A}$ kategori abelian. Buktikan bahwa
	$\cate{D}(\mathcal{A})$ merupakan kategori abelian jika dan hanya jika
	$\mathcal{A}$ terbelah.

% segment-id: o014.li2.chapter4.k-flat-derived-tensor-b.hn003
	\begin{hint}
		Untuk arah ``hanya jika'', benamkan $\mathcal{A}$ ke dalam
		$\cate{D}(\mathcal{A})$ dan terapkan hasil latihan-latihan sebelumnya.
		Untuk arah ``jika'', tunjukkan bagi $\mathcal{A}$ yang terbelah bahwa
		pengambilan kohomologi memberikan ekuivalensi dari
		$\cate{D}(\mathcal{A})$ ke $\mathcal{A}^{\Z}$.
	\end{hint}

% segment-id: o014.li2.chapter4.k-flat-derived-tensor-b.ex005
	\item Sering dikatakan bahwa aksioma oktahedral (TR5) bagi kategori
	bertriangulasi merupakan pengganti bagi teorema isomorfisme
	$\dfrac{X/Z}{Y/Z}\simeq X/Y$ dalam kategori abelian, dengan
	$X\supset Y\supset Z$ (Teorema~\ref{prop:Abel-cat-isom-thm} (ii)).
	Jelaskan secara eksplisit makna ungkapan tersebut.

% segment-id: o014.li2.chapter4.k-flat-derived-tensor-b.ex006
	\item ($\mathrm{K}_0$ kategori prabertriangulasi) Untuk setiap kategori
	prabertriangulasi $(\mathcal{D},T,\mathcal{H})$, definisikan
	$\mathrm{K}_0(\mathcal{D})$ sebagai hasil bagi modul-$\Z$ bebas yang
	dibangkitkan oleh $\Obj(\mathcal{D})$ terhadap relasi-relasi
% segment-id: o014.li2.chapter4.k-flat-derived-tensor-b.q009
	\[ \text{terdapat segitiga terbedakan}\; X \to Y \to Z \xrightarrow{+1} \implies [X] - [Y] + [Z] = 0, \]
	dengan $[X]\in\mathrm{K}_0(\mathcal{D})$ menyatakan kelas ekuivalensi yang
	memuat $X\in\Obj(\mathcal{D})$.
% segment-id: o014.li2.chapter4.k-flat-derived-tensor-b.l003
	\begin{enumerate}[(i)]
% segment-id: o014.li2.chapter4.k-flat-derived-tensor-b.i006
		\item Buktikan bahwa $[0]=0$, $[TX]=-[X]$, dan
		$[X\oplus Y]=[X]+[Y]$.
% segment-id: o014.li2.chapter4.k-flat-derived-tensor-b.i007
		\item Tunjukkan bahwa setiap funktor bertriangulasi
		$F:\mathcal{D}\to\mathcal{D}'$ menginduksi homomorfisme kanonik
		$\mathrm{K}_0(\mathcal{D})\to\mathrm{K}_0(\mathcal{D}')$.
% segment-id: o014.li2.chapter4.k-flat-derived-tensor-b.i008
		\item Tinjau funktor kohomologis $H:\mathcal{D}\to\mathcal{A}$, dengan
		$\mathcal{A}$ kategori abelian. Andaikan, untuk semua $X$,
		$|n|\gg_X0\implies H(T^nX)=0$. Definisikan homomorfisme alami
		$\mathrm{K}_0(\mathcal{D})\to\mathrm{K}_0(\mathcal{A})$.
		
% segment-id: o014.li2.chapter4.k-flat-derived-tensor-b.hn004
		\begin{hint}
			Petakan $[X]$ ke $\sum_n(-1)^n[H(T^nX)]$. Terapkan
			Lema~\ref{prop:K0-prep} (iv).
		\end{hint}
% segment-id: o014.li2.chapter4.k-flat-derived-tensor-b.i009
		\item Tinjau kasus khusus $\mathcal{D}:=\cate{D}^{\bdd}(\mathcal{A})$.
		Berikan homomorfisme yang saling invers
% segment-id: o014.li2.chapter4.k-flat-derived-tensor-b.g004
		\[\begin{tikzcd}[row sep=tiny]
			\mathrm{K}_0(\mathcal{A}) \arrow[shift left, r] & \mathrm{K}_0(\cate{D}^{\bdd}(\mathcal{A})) \arrow[shift left, l] \\
			{[A]} \arrow[mapsto, r] & {[A]} \\
			\sum_n (-1)^n {[\Hm^n(X)]} & {[X]} \arrow[mapsto, l]
		\end{tikzcd}\]
		dengan $A\in\Obj(\mathcal{A})$, yang juga dipandang sebagai kompleks
		yang terkonsentrasi pada derajat $0$.
% segment-id: o014.li2.chapter4.k-flat-derived-tensor-b.hn005
		\begin{hint}
			Satu komposisi jelas merupakan identitas; untuk komposisi lainnya,
			diperlukan pemenggalan kompleks.
		\end{hint}
	\end{enumerate}
	\index{grup K0@grup $\mathrm{K}_0$} \index[sym1]{K0}

% segment-id: o014.li2.chapter4.k-flat-derived-tensor-b.ex007
	\item Misalkan $X\in\Obj(\cate{D}(\mathcal{A}))$. Buktikan bahwa
% segment-id: o014.li2.chapter4.k-flat-derived-tensor-b.q010
	\begin{gather*}
		X \in \cate{D}^{\leq 0}(\mathcal{A}) \iff \forall Z \in \Obj(\cate{D}^{\geq 1}(\mathcal{A})), \; \Hom_{\cate{D}(\mathcal{A})}(X, Z) = 0, \\
		X \in \cate{D}^{\geq 0}(\mathcal{A}) \iff \forall Z \in \Obj(\cate{D}^{\leq -1}(\mathcal{A})), \; \Hom_{\cate{D}(\mathcal{A})}(Z, X) = 0.
	\end{gather*}
% segment-id: o014.li2.chapter4.k-flat-derived-tensor-b.hn006
	\begin{hint}
		Terapkan Proposisi~\ref{prop:derived-cat-orthogonality} (i), serta
		pasangan adjoin dan segitiga terbedakan dalam
		Proposisi~\ref{prop:derived-truncation-adjunction}.
	\end{hint}

% segment-id: o014.li2.chapter4.k-flat-derived-tensor-b.ex008
	\item Dalam situasi Teorema~\ref{prop:localization-triangulated-composite},
	andaikan $F'$ memenuhi
	$F'(\Obj(\mathcal{N}'))\subset\Obj(\mathcal{N}'')$ (anggap ini sebagai
	semacam ``keeksakan''). Jika $\mathcal{D}$ mempunyai subkategori
	$F$-injektif atau $F$-projektif, buktikan bahwa morfisme kanonik
	\[ \mathrm{R}^{\mathcal{N}''}_{\mathcal{N}} (F' F) \to \left( \mathrm{R}^{\mathcal{N}''}_{\mathcal{N}'} F' \right) \left( \mathrm{R}^{\mathcal{N}'}_{\mathcal{N}} F \right) \]
	atau
	\[ \left( \mathrm{L}^{\mathcal{N}''}_{\mathcal{N'}} F' \right) \left( \mathrm{L}^{\mathcal{N}'}_{\mathcal{N}} F \right) \to \mathrm{L}^{\mathcal{N}''}_{\mathcal{N}} (F' F) \]
	merupakan isomorfisme.
	
% segment-id: o014.li2.chapter4.k-flat-derived-tensor-b.ex009
	\item Dalam situasi Teorema~\ref{prop:localization-triangulated-composite},
	tunjukkan bahwa terdapat diagram komutatif berikut, yang terdiri atas
	funktor dan morfisme di antaranya (notasi $\mathcal{N}$ dan sebagainya
	dihilangkan):
% segment-id: o014.li2.chapter4.k-flat-derived-tensor-b.g005
	\[\begin{tikzcd}
		\mathrm{R}(F'' F' F) \arrow[r] \arrow[d] & \mathrm{R}(F'' F') \mathrm{R}F \arrow[d] \\
		\mathrm{R}F'' \mathrm{R}(F'F) \arrow[r] & \mathrm{R}F'' \mathrm{R}F' \mathrm{R}F
	\end{tikzcd} \quad \begin{tikzcd}
		\mathrm{L}(F'' F' F) & \mathrm{L}(F'' F') \mathrm{L}F \arrow[l] \\
		\mathrm{L}F'' \mathrm{L}(F'F) \arrow[u] & \mathrm{L}F'' \mathrm{L}F' \mathrm{L}F \arrow[u] \arrow[l] .
	\end{tikzcd}\]

% segment-id: o014.li2.chapter4.k-flat-derived-tensor-b.ex010
	\item Untuk $F=\Hom$, nyatakan secara eksplisit morfisme kanonik dari
	Proposisi~\ref{prop:bifunctor-cup} dalam bentuk
	$\Ext^{q-p}(X,Y)\to\Hom\left(\Hm^p(X),\Hm^q(Y)\right)$.

% segment-id: o014.li2.chapter4.k-flat-derived-tensor-b.ex011
	\item Rumuskan kembali bijeksi dalam Teorema~\ref{prop:Yoneda-Ext} dalam bentuk
	berikut. Dari ekstensi-$n$
	$0\to Y\to E^1\to\cdots\to E^n\to X\to0$, bangun
	$t^{-1}b\in\Hom_{\cate{D}(\mathcal{A})}(X,Y[n])$ seperti berikut:
% segment-id: o014.li2.chapter4.k-flat-derived-tensor-b.g006
	\[\begin{tikzcd}
		X \arrow[d, "b"']\arrow[phantom, r, "" {name=A, yshift=2ex}] & {} & & & X \arrow[d, "\identity"] \\
		W & E^1 \arrow[r] & E^2 \arrow[r] & \cdots \arrow[r] & X \\
		Y[n] \arrow[u, "{t: \text{kuasi-isomorfisme}}"] \arrow[phantom, r, "" {name=B, yshift=-2ex}] & Y \arrow[u] & & &
		\arrow[dash, from=A, to=B]
	\end{tikzcd}\]
	Dengan demikian diperoleh
	$t^{-1}b\in\Hom_{\cate{D}(\mathcal{A})}(X,Y[n])$. Buktikan bahwa bijeksi
	ini berbeda dengan faktor $(-1)^n$ dari bijeksi semula yang didasarkan
	pada $as^{-1}$. Untuk kasus khusus $n=1$, gunakan versi $t^{-1}b$ untuk
	menurunkan versi terkait dari Proposisi~\ref{prop:Ext1-via-ses} (yang
	melibatkan citra $\identity_Y$).
	
% segment-id: o014.li2.chapter4.k-flat-derived-tensor-b.hn007
	\begin{hint}
		Buktikan bahwa $ta+(-1)^{n-1}bs:Z\to W$ secara kanonik homotopik nol.
		Untuk bagian kedua, perhatikan tanda negatif dalam
		Proposisi~\ref{prop:ses-vs-triangle}.
	\end{hint}

% segment-id: o014.li2.chapter4.k-flat-derived-tensor-b.ex012
	\item Misalkan semua objek $S,T$ dalam kategori abelian $\mathcal{A}$
	memenuhi $\Ext^2_{\mathcal{A}}(S,T)=0$. Buktikan bahwa setiap
	$X\in\Obj(\cate{D}^{\bdd}(\mathcal{A}))$ isomorfik dengan
	$\bigoplus_n\Hm^n(X)[-n]$. Tunjukkan bahwa
	$\mathcal{A}:=\Z\dcate{Mod}$ memenuhi syarat ini.
	
% segment-id: o014.li2.chapter4.k-flat-derived-tensor-b.hn008
	\begin{hint}
		Kita boleh mengandaikan bahwa $X$ kompleks terbatas. Mulailah dari kasus
		$n\gg0$. Terapkan Teorema~\ref{prop:Yoneda-Ext} dan
		Proposisi~\ref{prop:Yoneda-zero} pada ekstensi-$2$
		$0\to\Ker(d^n)\to X^n\xrightarrow{d^n}X^{n+1}\to\Coker(d^n)\to0$,
		lalu ubah $X$ langkah demi langkah hingga $d^\bullet=0$, semuanya hingga
		isomorfisme.
	\end{hint}

% segment-id: o014.li2.chapter4.k-flat-derived-tensor-b.ex013
	\item Buktikan bahwa, jika $\mathcal{A}$ mempunyai cukup banyak objek
	injektif, $F:\mathcal{A}\to\mathcal{A}'$ funktor aditif eksak kiri, dan
	$\mathrm{R}F$ berdimensi hingga, maka homomorfisme
	$\mathrm{K}_0(\mathcal{A})\simeq\mathrm{K}_0(\cate{D}^{\bdd}(\mathcal{A}))
	\to\mathrm{K}_0(\cate{D}^{\bdd}(\mathcal{A}'))\simeq
	\mathrm{K}_0(\mathcal{A}')$ yang diinduksi oleh funktor bertriangulasi
	$\mathrm{R}F$ ditentukan oleh
	\[ [A] \mapsto \sum_n (-1)^n \left[\mathrm{R}^n F(A)\right], \quad A \in \Obj(\mathcal{A}). \]

% segment-id: o014.li2.chapter4.k-flat-derived-tensor-b.ex014
	\item Misalkan $R$ gelanggang. Buktikan bahwa limit langsung terfilter
	$\varinjlim$ dari kompleks-kompleks K-datar dalam
	$\cate{C}(R\dcate{Mod})$ tetap K-datar.
	
% segment-id: o014.li2.chapter4.k-flat-derived-tensor-b.ex015
	\item Misalkan $R$ gelanggang. Jika kompleks modul-$R$ $P$ K-datar dan
	setiap $P^n$ merupakan modul-$R$ datar, maka $P$ disebut
	\emph{K-datar kuat}.
% segment-id: o014.li2.chapter4.k-flat-derived-tensor-b.l004
	\begin{enumerate}[(i)]
% segment-id: o014.li2.chapter4.k-flat-derived-tensor-b.i010
		\item Buktikan bahwa, untuk setiap
		$X\in\Obj(\cate{C}(R\dcate{Mod}))$, terdapat kompleks K-datar kuat $P$
		beserta kuasi-isomorfisme $P\to X$.
% segment-id: o014.li2.chapter4.k-flat-derived-tensor-b.hn009
		\begin{hint}
			Gunakan dual argumen keberadaan resolusi K-injektif; lihat
			\S\ref{sec:K-injectives}.
		\end{hint}
% segment-id: o014.li2.chapter4.k-flat-derived-tensor-b.i011
		\item Tuliskan $\mathcal{SF}$ untuk subkategori penuh
		$\cate{K}(R\dcate{Mod})$ yang dibentuk oleh kompleks-kompleks K-datar
		kuat. Nyatakan kategori $\cate{D}(R\dcate{Mod})$ sebagai lokalisasi
		$\mathcal{SF}$ terhadap kuasi-isomorfisme.
	\end{enumerate}

% segment-id: o014.li2.chapter4.k-flat-derived-tensor-b.ex016
	\item (P.\ Berthelot, A.\ Ogus) Misalkan $f$ bukan pembagi nol dalam
	gelanggang komutatif $R$. Untuk setiap modul-$R$ $M$, tuliskan
	$M[f]:=\{m\in M:fm=0\}$. Modul $M$ dengan $M[f]=0$ disebut
	\emph{bebas torsi-$f$}. Pertama, tunjukkan bahwa $M$ yang bebas torsi-$f$
	dapat dibenamkan sebagai submodul dari
	$M[\frac{1}{f}]=M\dotimes{R}R[\frac{1}{f}]$. Karena itu, $f^mM$
	bermakna untuk setiap $m\in\Z$.

	Misalkan $X$ kompleks yang tersusun atas modul-$R$ bebas torsi-$f$.
	Berdasarkan pengamatan di atas, definisikan subkompleks $\eta_fX$ dari
	$X^\bullet\dotimes{R}R[\frac{1}{f}]$ dengan
	\[ (\eta_f X)^n := \left\{ x \in f^n X^n : d_X^n(x) \in f^{n+1} X^{n+1} \right\}, \quad n \in \Z . \]
% segment-id: o014.li2.chapter4.k-flat-derived-tensor-b.l005
	\begin{enumerate}[(i)]
% segment-id: o014.li2.chapter4.k-flat-derived-tensor-b.i012
		\item Buktikan bahwa perkalian suku demi suku dengan $f$ memberikan
		isomorfisme $\eta_f(X[1])\rightiso(\eta_fX)[1]$.
% segment-id: o014.li2.chapter4.k-flat-derived-tensor-b.i013
		\item Tunjukkan bahwa, jika $f,g\in R$ keduanya bukan pembagi nol dan
		setiap $X^n$ bebas torsi-$fg$, maka
		$\eta_f\eta_gX=\eta_{fg}X$.
% segment-id: o014.li2.chapter4.k-flat-derived-tensor-b.i014
		\item Berikan isomorfisme kanonik
		$\Hm^n(X)/\Hm^n(X)[f]\rightiso\Hm^n(\eta_fX)$; jadi efek $\eta_f$
		ialah mengoreksi torsi-$f$. Dari sini, simpulkan bahwa jika $X$ dan $Y$
		merupakan kompleks yang tersusun atas modul-$R$ bebas torsi-$f$ dan
		$\alpha:X\to Y$ suatu kuasi-isomorfisme, maka $\alpha$ menginduksi
		kuasi-isomorfisme $\eta_f(\alpha):\eta_fX\to\eta_fY$.
% segment-id: o014.li2.chapter4.k-flat-derived-tensor-b.i015
		\item Buktikan bahwa $\eta_f$ menginduksi funktor aditif
		$\mathrm{L}\eta_f:\cate{D}(R\dcate{Mod})\to\cate{D}(R\dcate{Mod})$
		dengan cara ini, yang pembatasannya ke $\mathcal{SF}$ dari latihan
		sebelumnya sama dengan $\eta_f$.
% segment-id: o014.li2.chapter4.k-flat-derived-tensor-b.i016
		\item Tunjukkan bahwa $\mathrm{L}\eta_f$ bukan funktor
		bertriangulasi. Jadi ia bukan funktor turunan kiri dalam arti
		Definisi~\ref{def:Verdier-localization-functor}, meskipun tetap merupakan
		ekstensi Kan kanan.

% segment-id: o014.li2.chapter4.k-flat-derived-tensor-b.hn010
		\begin{hint}
			Ambil $R=\Z$ dan $f=p$ prima. Verifikasi bahwa
			$\mathrm{L}\eta_p(\Z/p\Z)=0$, sedangkan
			$\mathrm{L}\eta_p(\Z/p^2\Z)\simeq\Z/p\Z$.
		\end{hint}
	\end{enumerate}

% segment-id: o014.li2.chapter4.k-flat-derived-tensor-b.ex017
	\item Melanjutkan latihan sebelumnya, berikan morfisme kanonik
	$\mathrm{L}\eta_f X \otimesL[R] \mathrm{L}\eta_f Y \to
	\mathrm{L}\eta_f\left( X \otimesL[R] Y \right)$ dan tunjukkan bahwa
	morfisme ini simetris terhadap $X$ dan $Y$ dalam arti yang sesuai.
	Sesungguhnya, dalam bahasa \cite[Catatan 3.1.8]{Li1},
	$\mathrm{L}\eta_f$ merupakan funktor monoidal longgar kanan terhadap
	$\otimesL[R]$.
\end{Exercises}
